\TieudeT{PHONG TỎA VẬT LÍ 11 - DAO ĐỘNG VÀ SÓNG}
\subsubsection*{PHẦN I. Thí sinh trả lời câu hỏi từ câu 1 đến câu 18. Mỗi câu hỏi thí sinh chỉ chọn một phương án}
\setcounter{ex}{0}
	\begin{ex} %[3P1Y1-1][Loai1] 
		Một vật dao động điều hòa theo phương trình $x = A\cos\left(\omega t + \varphi \right)$ ($\omega  > 0$). Tần số góc của dao động là
		\choice
		{$A$}
		{\True $\omega$}
		{$\varphi$}
		{$x$}
		\loigiai{
			Tần số góc của dao động là $\omega$.
		}
		%\haidong{5}
	\end{ex}
	%\loai{Xác định biên độ dao động và tần số góc}
	\begin{ex}%[3P1Y1-2][Loai1]
		Một chất điểm dao động theo phương trình $x = 6\cos \omega t\ cm$. Dao động của chất điểm có biên độ là
		\choice
		{\True $6\ cm$}
		{$3\ cm$}
		{$2\ cm$}
		{$12\ cm$}
		\loigiai{
			Phương trình dao động tổng hợp là: $x = A\cos \omega t\ cm$ $\Rightarrow$ $A = 6\ cm$.
		}
		
		%\haidong{5}
	\end{ex}
	%\loai{Lí thuyết vận tốc}
	\begin{ex}%[3P1Y2-1][Loai1]
		Trong quá trình dao động, vận tốc của vật có giá trị lớn nhất (cực đại) khi vật
		\choice
		{đi qua vị trí cân bằng}
		{\True đi qua vị trí cân bằng theo chiều dương}
		{đi qua vị trí cân bằng theo chiều âm}
		{ở biên}
		\loigiai{
			Trong quá trình dao động, vận tốc của vật có giá trị lớn nhất (cực đại) khi vật đi qua vị trí cân bằng theo chiều dương.
		}
		%\haidong{5}
	\end{ex}
	%\loai{Dùng công thức độc lập x, v tìm x}
	\begin{ex}%[3P1B2-2][Loai1]
		Một vật dao động điều hòa với biên độ A, vận tốc cực đại $v_{\max}$. Vật có tốc độ $0,6v_{\max}$ khi li độ của vật có độ lớn là
		\choice
		{\True $0,8A$}
		{$0,6A$}
		{$0,4A$}
		{$0,5A$}
		\loigiai{
			\begin{itemize}
				\item Luôn có: $\left(\dfrac{x}{A}\right)^2+\left(\dfrac{v}{v_{\max}}\right)^2=1\quad  (*)$
				\item Khi $\left| v \right|=0,6v_{\max}\quad (*) \Rightarrow \left(\dfrac{x}{A}\right)^2+\left(\dfrac{0,6v_{\max}}{v_{\max}}\right)^2=1\Rightarrow \left|x\right|=0,8A$.
			\end{itemize}
		}
		
		
		%\haidong{14}
	\end{ex}
	\begin{ex} %[3P1B4-1][Loai1]  
		Biết pha ban đầu của một vật dao động điều hòa, ta xác định được
		\choice
		{quỹ đạo dao động}
		{cách kích thích dao động}
		{chu kì và trạng thái dao động}
		{\True chiều chuyển động của vật lúc ban đầu}
		\loigiai{
			Pha dao động cho ta biết trạng thái chuyển động của vật.
		}
		%\haidong{5}
	\end{ex}
	%\loai{Khoảng thời gian ngắn nhất từ li độ $x_1$ đến li độ $x_2$ vị trí đẹp}
	\begin{ex}%[3P1B8-1][Loai1]
		Một vật dao động điều hòa với biên độ A, tần số 5 Hz. Thời gian ngắn nhất để vật đi từ vị trí có li độ $x_1 = - 0,5A$ đến vị trí có li độ $x_2 = 0,5A$ là
		\choice
		{$1\ s$}
		{$\dfrac{1}{10}\ s$}
		{\True $\dfrac{1}{30}\ s$}
		{$\dfrac{1}{20}\ s$}
		\loigiai{
			\begin{center}
				\begin{tikzpicture}
					\tkzDefPoints{0/0/o, 5/0/a, 4.325/0/b, 3.525/0/c, 2.5/0/d,  -5/0/h, -4.325/0/g, -3.525/0/f, -2.5/0/e, -6/0/x', 6/0/x}
					\tkzDefShiftPoint[e](-90:0.5cm){e1}
					\tkzDefShiftPoint[o](-90:0.5cm){o1}
					\tkzDefShiftPoint[d](-90:0.5cm){d1}
					\tkzDrawSegments[blue,->,very thick](x',x)
					\tkzDrawPoints[size=5,fill=black](o,a,b,c,d,e,f,g,h)
					\tkzLabelPoint[above](o){\tiny $O$}
					\tkzLabelPoint[above](a){\tiny$A$}
					\tkzLabelPoint[above](b){\tiny$\dfrac{A\sqrt{3}}{2}$}
					\tkzLabelPoint[above](c){\tiny$\dfrac{A\sqrt{2}}{2}$}
					\tkzLabelPoint[above](d){\tiny$\dfrac{A}{2}$}
					\tkzLabelPoint[above](h){\tiny$-A$}
					\tkzLabelPoint[above](g){\tiny $-\dfrac{A\sqrt{3}}{2}$}
					\tkzLabelPoint[above](f){\tiny$-\dfrac{A\sqrt{2}}{2}$}
					\tkzLabelPoint[above](e){\tiny$-\dfrac{A}{2}$}
					\tkzLabelPoint[above](x){\tiny$x$}
					\tkzDrawSegments[->,very thick,red](e1,o1 o1,d1)
					\tkzLabelSegment[below](e1,o1){\tiny$\Delta t_1=\dfrac{T}{12}$}
					\tkzLabelSegment[below](o1,d1){\tiny$\Delta t_2=\dfrac{T}{12}$}
				\end{tikzpicture}
			\end{center}
			Ta có: $\Delta t=\Delta t_1+\Delta t_2=\dfrac{T}{6}=\dfrac{1}{30}\ s$.
		}
		
		%\haidong{20}
	\end{ex}
	%\loai{Khoảng thời gian li độ lớn hơn (hoặc nhỏ hơn) b trong một chu kì}
	\begin{ex}%[3P1K8-1][Loai1]
		Vật dao động điều hoà với biên độ $6\ cm$, chu kì $1,2\ s$. Trong một chu kì, khoảng thời gian để li độ ở trong khoảng $(- 3\ cm;\ 3\ cm)$ là
		\choice
		{$0,3\ s$}
		{$0,2\ s$}
		{$0,6\ s$}
		{\True $0,4\ s$}
		\loigiai{
			\begin{itemize}
				\item Ta thấy $[-3\ cm;\ 3\ cm] = [-\dfrac{A}{2};\ \dfrac{A}{2}]$.
				\item Trong nửa chu kì thì vật đi từ vị trí có li độ $\dfrac{A}{2}$ đến $-\dfrac{A}{2}$ sẽ ứng với góc quay được là $\varphi = \dfrac{\pi}{3}\ rad$ và hết khoảng thời gian $t_1 = \dfrac{T}{6}$. 
				\item Trong nửa chu kì còn lại vật sẽ có li độ $[-3\ cm;\ 3\ cm]$ trong khoảng thời gian $t_1$ nhưng theo chiều ngược lại.
				\item Trong một chu kì, khoảng thời gian để li độ của vật trong khoảng $[-3\ cm;\ 3\ cm]$ là $\Delta t = 2t_1 = \dfrac{T}{3} = 0,4\ s$.
			\end{itemize}
		}
		
		%\haidong{20}
	\end{ex}
	%\loai{Thời điểm vật đạt tốc độ cho trước lần thứ n}
	\begin{ex}%[3P1K9-1][Loai1]  
		Một vật dao động với phương trình $x=6\cos\left( \dfrac{10\pi t}{3} \right)\ cm$. Tính từ $t = 0$ thời điểm lần thứ $2013$ vật có tốc độ $10\pi \ cm/s$ là
		\choice
		{$302,35\ s$}
		{\True $301,85\ s$}
		{$302,05\ s$}
		{$302,15\ s$}
		\loigiai{
			\begin{itemize}
				\item Chu kì: $T=\dfrac{2\pi }{\omega }=0,6\ s$.
				\item Thay tốc độ $10\pi \ cm/s$ vào phương trình: $x^2+\dfrac{v^2}{\omega ^2}=A^2\Rightarrow \left| x\right|=3\sqrt{3}\ cm$. 
				\item Ta nhận thấy: $\dfrac{2013}{4}=503$ dư 1 
				$\Rightarrow t=503T+t_1$ nên ta chỉ cần tìm $t_1$:\\
				$t_1=\dfrac{T}{12}\Rightarrow t=503T+\dfrac{T}{12}=301,85\ s$. 
		\end{itemize}}
		
		
		%\haidong{26}
	\end{ex}
	%\loai{Hai thời điểm ngược pha - vận tốc và gia tốc}
	\begin{ex}%[3P1KC-1][Loai1]
		Một vật dao động điều hoà trên trục Ox, tại thời điểm t nào đó vận tốc có giá trị âm và gia tốc có giá trị dương. Tại thời điểm $t + \dfrac{T}{2}$ thì
		\choice
		{vận tốc và gia tốc có giá trị âm}
		{\True vận tốc có giá trị dương, gia tốc có giá trị âm}
		{vận tốc và gia tốc có giá trị dương}
		{vận tốc có giá trị âm, gia tốc có giá trị dương}
		\loigiai{
			Thời điểm t, vật có $v < 0,\ a > 0$, tức vật đang có li độ âm (li độ và gia tốc luôn trái dấu) và đi theo chiều âm, sau $\dfrac{T}{2}$ trạng thái dao động ngược lại nên vật có $v > 0, a < 0$ (li độ dương).
		}
		
		%\haidong{22}
	\end{ex}
	%\loai{Tốc độ trung bình trong khoảng thời gian từ $t_1$ đến $t_2$} 
	\begin{ex}%[3P1KD-3][Loai1]
		Cho chất điểm M đang chuyển động tròn đều trên một quỹ đạo tròn có bán kính bằng $30\ cm$, với tốc độ góc là $\pi\ rad/s$. Gọi P là hình chiếu của điểm M xuống một đường thẳng đi qua tâm và nằm trong mặt phẳng quỹ đạo. Tốc độ chuyển động trung bình của điểm P trong quãng thời gian $3\ s$ là
		\choice
		{$20\ cm/s$}
		{\True $60\ cm/s$}
		{$\dfrac{20}{3}\ cm/s$}
		{$40\ cm/s$}
		\loigiai{
			\begin{itemize}
				\item Trong khoảng thời gian 3 s chất điểm quay được góc: $\Delta \varphi = 3\pi = 2\pi + \pi\Rightarrow  \Delta s = 4A+2A= 180\ cm$.
				\item Tốc độ trung bình bằng: $180:3 = 60\ cm/s$.
			\end{itemize}
		}
		%\haidong{27}
	\end{ex}
	\begin{ex}%[3P1-19.2K][Loai1] 
		\immini[thm]{
			Cho đồ thị hai dao động điều hòa như hình vẽ. Độ lệch pha của chúng là
			\choice
			{$\pi$}
			{$3\pi$}
			{$\dfrac{\pi}{2}$}
			{\True $0$}}
		{
			\begin{tikzpicture}[draw=Mapcolor,scale=1,thick,>=stealth']
				\pgfplotsset{width=8cm,height=4cm,compat=1.4}
				\begin{axis}[
					axis lines=middle,% Tùy chọn hệ trục tọa độ Đề các
					%clip=true,
					xmin=0,xmax=5.5,xstep=1,ymin=-2.16,ymax=2.3,ystep=1,
					grid=major, %Có các tùy chọn: minor|major|both|none
					x grid style={dashed,Mapcolor},
					y grid style={dashed,Mapcolor},
					ytick={-2,0,1.2,2},
					xtick={1,2,3,4,5},
					xticklabels={,0.01,,0.02},
					yticklabels={$-A_1$,$0$,$A_2$,$A_1$},
					xticklabel style={black,below=6pt},
					xlabel style={right=3pt},
					xlabel=$t(s)$,
					ylabel style={above=3pt},
					ylabel={$x\ (cm)$}
					]
					\addplot[line width=1.2pt,domain=0:5,samples=200,Mapcolor]{2*cos(deg(0.5*pi*x))};
					\addplot[line width=1.2pt,domain=0:5,samples=200,Mapcolor,dashed]{1.2*cos(deg(0.5*pi*x))};
				\end{axis}
			\end{tikzpicture}
		}
		\loigiai{
			Hai dao động cùng pha $\Delta \varphi =2k\pi $}
		
		%\haidong{12}
	\end{ex}
	%\loai{Biểu thức}
	\begin{ex} %[3P1Y5-1][Loai1] 
		Một con lắc lò xo có khối lượng vật nhỏ là m dao động điều hòa theo phương ngang với phương trình $x = A\cos\omega t$. Mốc tính thế năng ở vị trí cân bằng. Cơ năng của con lắc là
		\choice
		{$m\omega A^2$}
		{$\dfrac{1}{2}m\omega A^{2}$}
		{$m\omega ^{2}A^{2}$}
		{\True $\dfrac{1}{2}m\omega ^{2}A^{2}$}
		\loigiai{
			Cơ năng của vật dao động điều hòa: $W = W_{\text{đ}\max} = \dfrac{1}{2}mv_{\max}^{2} = \dfrac{1}{2}m\omega ^2A^2$.
		}
		%\haidong{6}
	\end{ex}
	%\loai{Tính tần số góc}
	\begin{ex} %[3P1Y3-2][Loai1] 
		Một con lắc lò xo có $k = 40\ N/m$ và $m = 100\ g$. Dao động riêng của con lắc này có tần số góc là
		\choice
		{$400\ rad/s$}
		{$0,1\pi\ rad/s$}
		{\True $20\ rad/s$}
		{$0,2\pi\ rad/s$}
		\loigiai{
			Tần số góc: $\omega =\sqrt{\dfrac{k}{m}}=20\ rad/s$.
		}
		
		%\haidong{17}
	\end{ex}
	\begin{ex}%[3P1-20.1K][Loai1]
		\immini[thm]{Một vật nhỏ dao động điều hòa với tần số góc $10\ rad/s$. Giá trị còn thiếu trong dấu ? ở đồ thị hình bên là
			\motcot
			{$400$}
			{\True $-4$}
			{$40$}
			{$-400$}}{\begin{tikzpicture}[Mapcolor]
				\draw[->](0,-1.5)--(0,1.5)node[left]{$a\ (m/s^2)$};
				\draw[->](-2.5,0)--(2.5,0)node[below right]{$x\ (cm)$};
				\draw[very thick,Mapcolor](-2,1)--(2,-1);
				\draw[dashed](-2,0)node[below]{$-4$}--(-2,1)--(0,1);
				\draw[dashed](2,0)node[above]{$4$}--(2,-1)--(0,-1)node[left]{$?$};
		\end{tikzpicture}}
		\loigiai{
			Ta có: $a_{\max} = \omega^2A = 400\ cm/s^2 = 4\ m/s^2$.
		}
		
		%\haidong{9}
	\end{ex}
	
	
	
	%\loai{Công thức}
	\begin{ex}
		Một con lắc đơn đang dao động điều hòa với phương trình $s = S_0\cos(\omega t + \varphi)\ (S_0> 0)$. Đại lượng $S_0$ được gọi là
		\choice
		{\True biên độ dài của dao động}			
		{tần số của dao động}
		{li độ góc của dao động}
		{pha ban đầu của dao động}
		%\haidong{5}
	\end{ex}
	%\loai{Tìm độ biến dạng tại vị trí cân bằng }
	\begin{ex}%[3P1B6-2][Loai1]%Bài 15 dạng 1
		Một chất điểm $M$ chuyển động tròn đều trên quỹ đạo tâm $O$ bán kính $8\ cm$ với tốc độ dài $200\ cm/s$. Gọi $P$ là hình chiếu của $M$ lên trục $Ox$ nằm trong mặt phẳng quỹ đạo. Khi $P$ cách $O$ một đoan $b$ thì nó có tốc độ là $100 \sqrt{2} \ cm/s$. Giá trị của $b$ là
		\choice
		{$4\ cm$}
		{\True $4 \sqrt{2} \ cm$}
		{$4 \sqrt{3} \ cm$}
		{$2\ cm$}
		\loigiai{
		}
		
		%\haidong{20}
	\end{ex}
	%\loai{Lí thuyết vận tốc}
	\begin{ex}%[3P1-16.1B][Loai1]  
		Vận tốc của con lắc đơn có vật nặng khối lượng m, chiều dài dây treo $\ell$, dao động với biên độ góc $\alpha _m$ khi qua li độ góc $\alpha $ được tính theo biểu thức nào sau đây?
		\choice
		{$v^2=2mg\ell\left({\cos\alpha - \cos\alpha _m}\right)$}
		{$v^2=mg\ell\left({\cos\alpha _m-\cos\alpha }\right)$}
		{\True $v^2=2g\ell\left({\cos\alpha -\cos\alpha _m}\right)$}
		{$v^2=mg\ell\left({\cos\alpha - \cos\alpha _m}\right)$}
		\loigiai{
			Vận tốc của con lắc đơn $v^2=2g\ell\left({\cos \alpha -\cos \alpha _m}\right)$.
		}
		%\haidong{10}
	\end{ex}
	
	
	%\loai{Công thức độc lập}
	\begin{ex} %[3P1-16.3][Loai1]  
		Một con lắc đơn có chiều dài $\ell = 1$ m được kéo ra khỏi vị trí cân bằng một góc $\alpha _0=5^\circ$ so với phương thẳng đứng rồi thả nhẹ cho vật dao động. Cho $g=\pi ^2=10\ m/s^2$. Tốc độ của con lắc khi về đến giá trị cân bằng có giá trị là
		\choice
		{$15,8\ m/s$}
		{\True $0,278\ m/s$}
		{$0,028\ m/s$}
		{$0,087\ m/s$}
		\loigiai{
			Vận tốc con lắc khi qua vị trí cân bằng:  $v_{\max}=\omega.S_0=\sqrt{\dfrac{g}{\ell}}.\ell.\alpha_0=0,278$ m/s.
		}
		
		%\haidong{28}
	\end{ex}
	

\subsubsection*{PHẦN 2. Thí sinh trả lời từ câu 1 đến câu 4. Trong mỗi ý a), b), c), d) ở mỗi câu, thí sinh chọn đúng hoặc sai.}
\setcounter{ex}{0}
\begin{ex}
	Một vật dao động điều hòa với phương trình: $x=10 \cos \left(2 \pi t-\dfrac{\pi}{2}\right)\ cm$ (t đo bằng giây).
	\choiceTF[t]
	{Thời gian ngắn nhất vật đi từ vị trí có li độ $-10 \ cm$ đến vị trí có li độ $-5 \ cm$ là $0,5 \ s$}
	{Thời gian ngắn nhất vật đi từ vị trí có li độ $5 \sqrt{3} \ cm$ đến vị trí cách vị trí cân bằng $5 \ cm$ là $\dfrac{1}{4} \ s$}
	{\True Thời gian ngắn nhất vật đi từ vị trí cân bằng đến vị trí có li độ $5 \sqrt{2} \ cm$ là $\dfrac{1}{8}\ s$}
	{\True Thời gian ngắn nhất vật đi từ vị trí cách vị trí cân bằng $5 \ cm$ đến vị trí có li độ $-5 \sqrt{2} \ cm$ là $\dfrac{1}{24} \ s$}
	\loigiai{
		\begin{itemchoice}
			\itemch Sai.
			\itemch Sai.
			\itemch Đúng.
			\itemch Đúng.
		\end{\itemchoice}
	}
	
	%\haidong{25}
\end{ex}
\begin{ex}
	Cho một vật dao động điều hòa với phương trình $x=A\cos(\omega t+\varphi)$.
	\choiceTF[t]
	{Khoảng thời gian ngắn nhất sau đó trạng thái dao động lặp lại như cũ gọi là tần số dao động}
	{\True Đại lượng không đổi theo thời gian là biên độ dao động}
	{Quỹ đạo dao động của vật được biểu diễn dưới dạng hình sin}
	{\True Pha ban đầu của dao động phụ thuộc vào gốc thời gian và chiều dương của hệ tọa độ}
	\loigiai{
		\begin{itemchoice}
			\itemch Sai.
			\itemch Đúng.
			\itemch Sai.
			\itemch Đúng.
		\end{\itemchoice}
	}
	%\haidong{25}
\end{ex} 

\begin{ex}
	Một chất điểm dao động điều hòa với biên độ $6 \ cm$ và chu kì $2 \ s$.
	\choiceTF[t]
	{Thời gian dài nhất để chất điểm đi được quãng đường bằng $18 \ cm$ là $\dfrac{4}{3} \ s$}
	{\True Thời gian ngắn nhất để chất điểm đi được quãng đường bằng $6 \sqrt{3} \ cm$ là $\dfrac{2}{3} \ s$}
	{Thời gian dài nhất để chất điểm đi được quãng đường $6 \ cm$ là $0,4 \ s$}
	{Thời gian dài nhất để chất điểm đi quãng được đường bằng $66 \ cm$ là $\dfrac{16}{3} \ s$}
	\loigiai{
		\begin{itemchoice}
			\itemch Sai.
			\itemch Đúng.
			\itemch Sai.
			\itemch Sai
		\end{\itemchoice}
	}
	
	%\haidong{25}
\end{ex}
\begin{ex}
	Một con lắc lò xo nằm ngang dao động điều hòa  với lò xo có độ cứng $k$, vật nặng có khối lượng $m$.
	\choiceTF[t]
	{\True Vị trí cân bằng của dao động cũng là vị trí lò xo có chiều dài tự nhiên}
	{\True Chu kì dao động là $T=2 \pi \sqrt{\dfrac{m}{k}}$}
	{Tần số dao động là $f=\dfrac{1}{2 \pi} \sqrt{\dfrac{m}{k}}$}
	{\True Khoảng thời gian vật đi từ biên về vị trí cân bằng là $t=\dfrac{T}{4}$}
	\loigiai{
		\begin{itemchoice}
			\itemch Đúng.
			\itemch Đúng.
			\itemch Sai. Tần số dao động là $f=\dfrac{1}{2 \pi} \sqrt{\dfrac{k}{m}}$.
			\itemch Đúng.
		\end{\itemchoice}
	}
	
	%\haidong{25}
\end{ex}
\subsubsection*{PHẦN 3. Thí sinh trả lời từ câu 1 đến câu 6.}
\setcounter{ex}{0}
\begin{ex}
	Một chất điểm dao động điều hoà với phương trình $x=8 \cos 5 t$ (x tính bằng $cm$, $t$ tính bằng $s$). Tốc độ của chất điểm khi đi qua vị trí cân bằng là bao nhiêu $cm/s$?
	\shortans[oly]{40}
	\loigiai{
		$40 \ cm/s$
	}
\end{ex}
\begin{ex}
	Một chất điểm dao động điều hòa theo phương trình $x=5 \sin \left(2 \pi t+\dfrac{\pi}{3}\right)\ cm$. Thời điểm chất điểm đi qua vị trí có li độ $-2,5 \ cm$ theo chiều dương lần thứ $2024$ là bao nhiêu phút kể từ thời điểm ban đầu (làm tròn kết quả đến chữ số hàng phần mười)?
	\shortans[oly]{33,7} \loigiai{
		33,7 phút
	}
	
	%\haidong{30}
\end{ex}
\begin{ex}
	Một vật dao động điều hòa với chu kì $T$. Thời gian ngắn nhất vật đi từ vị trí gia tốc vật có giá trị cực đại đến vị trí vận tốc vật có giá trị cực tiểu là bao nhiêu chu kì?
	\shortans[oly]{0,75}
	\loigiai{
		0,75
	}
	
	%\haidong{30}
\end{ex}
\begin{ex}
	Một chất điểm dao động điều hòa với phương trình $x=4 \cos \left(4 \pi t+\dfrac{2 \pi}{3}\right)\ (cm)$ ($t$ tính bằng $s$). Quãng đường lớn nhất chất điểm đi được bằng bao nhiêu $cm$ trong khoảng thời gian $\Delta t=\dfrac{1}{8} \ s$ (làm tròn kết quả đến chữ số hàng phần trăm)?
	\shortans[oly]{5,66}
	\loigiai{
		$5,66 \ cm$
	}
\end{ex}
\begin{ex}
	Con lắc lò xo treo thẳng đứng, dao động điều hòa với phương trình $x=5 \cos \left(4 \pi t+\dfrac{\pi}{3}\right)\ cm$. Chiều dài tự nhiên của lò xo là $40 \ cm$. Chọn chiều dương hướng lên. Lấy $g=\pi^2 \ m/s^2$. Chiều dài của lò xo là bao nhiêu  centimet khi vật dao động được $\dfrac{2T}{3}$, kể từ thời điểm $t=0$ (làm tròn kết quả đến chữ số hàng phần mười)?
	\shortans[oly]{43,8}
	\loigiai{43,8}
\end{ex}
%\loai{Tìm tốc độ tại vị trí bất kì}
\begin{ex}%[3P1K6-2][Loai1]%Bài 15 dạng 1
	Một con lắc lò xo treo thẳng đứng tại nơi có $g = 10\ m/s^2$. Vật đang cân bằng thì lò xo giãn $5\ cm$. Kéo vật xuống dưới vị trí cân bằng $1\ cm$ rồi truyền cho nó tốc độ $v_0$. Sau đó vật dao động điều hòa với tốc độ cực đại $30\sqrt{2}\ cm/s$. Giá trị của $v_0$ là bao nhiêu cm/s?
	\shortans[oly]{40}
	\loigiai{
		\begin{itemize}
			\item Ta có: $\Delta \ell_0 = 5\ cm \Rightarrow \omega  = 10\sqrt{2}\  rad/s,\ v_{\max} =  30\sqrt{2}\ cm/s \Rightarrow A = 3\ cm$.
			\item Khi truyền tốc độ $v_0$, vật đang có $\left|x\right| = 1\ cm \Rightarrow v_0=\omega \sqrt{A^2-x^2}=  40\ cm/s.$	
		\end{itemize}
	}
	
	%\haidong{30}
\end{ex}
